\section{Beurteilung ausgew\"ahlter Handlungsoptionen}
\label{sec:optionen}

Ausgangslage laut Aufgabenstellung und Gesch\"aftsbericht 2024: Die Aumann~AG
verf\"ugt \"uber eine Nettofinanzliquidit\"at von \MioEUR{\NettoliquiditaetZFIV}
zum Jahresende 2024\quelleGB{2024}{5} und leitet daraus die M\"oglichkeit ab,
sowohl organisch als auch \"uber gezielte Akquisitionen zu wachsen. Zugleich
ging der Auftragseingang 2024 um \PctBetrag{\AuftragseingangDeltaZFIV} auf
\MioEUR{\AuftragseingangZFIV} zur\"uck.\quelleGB{2024}{2}

Die drei nachfolgend beurteilten Sto{\ss}richtungen sind die in
Abschnitt~\ref{sec:strategie} sachlich dargestellten Ma{\ss}nahmen. Ihre
quantitative Unterlegung (Investitionssummen, Kapitalbindung und erwarteter
EBITDA-Zuwachs) folgt den Vorgaben der Aufgabenstellung und wird in
Abschnitt~\ref{sec:cashflow} durchgerechnet; die dortigen Tabellen sind die
Zahlenbasis der hier vorgenommenen Einsch\"atzungen.

\subsection{Diversifikation in Zukunftsbranchen (Option A)}
\label{subsec:option-a}

\subsubsection{Relevanz der Strategie (Management-Summary)}
% FAKTENBASIS 5.1.1
%  - Auftragseingang 2024: \AuftragseingangZFIV Mio. EUR, Rueckgang
%    \AuftragseingangDeltaZFIV % (GB 2024, S. 2)
%  - Nettofinanzliquiditaet 31.12.2024: \NettoliquiditaetZFIV Mio. EUR (GB 2024, S. 5)
%  - Anteil E-Mobility am Geschaeft: \AnteilEMobility % (GB 2024, S. 10)
%    -> Klumpenrisiko gegenueber dem Investitionszyklus der Automobilindustrie
%  - Segment Next Automation: Umsatz \UmsatzNextAutoZFIII -> \UmsatzNextAutoZFIV
%    -> \UmsatzNextAutoZFV Mio. EUR, Marge \EbitdaMargeNextAutoZFIII ->
%    \EbitdaMargeNextAutoZFIV -> \EbitdaMargeNextAutoZFV %
%  - Auftragseingang Next Automation: \AuftragseingangNextAutoZFIV ->
%    \AuftragseingangNextAutoZFV Mio. EUR (Tabelle \ref{tab:next-automation})
%  - Zielmaerkte laut GB 2024, S. 10: Clean Tech, Aerospace, Life Sciences
%  - Treiber laut Aufgabenstellung: Automatisierungs- und Robotikbedarf,
%    Reshoring, demografischer Wandel
%  - Investitionsannahme Aufgabe 6: 35 Mio. EUR ueber drei Jahre,
%    +16 Mio. EUR EBITDA ab 2028
\bk{Aumann erwirtschaftet mit rund \Pct{\AnteilEMobility} den weit
\"uberwiegenden Teil des Gesch\"afts mit Produktionsanlagen f\"ur die
Elektromobilit\"at. Die Auslastung des Unternehmens h\"angt damit an den
Investitionsentscheidungen einer einzigen Kundenbranche. Wie unmittelbar
das wirkt, zeigt das Jahr 2024: Umsatz und Ergebnis erreichten H\"ochstwerte,
der Auftragseingang fiel gleichwohl um
\PctBetrag{\AuftragseingangDeltaZFIV}, und der Aktienkurs folgte dem
Auftragseingang, nicht dem Ergebnis (Abschnitt~\ref{sec:aktienkurs}).
Option~A setzt genau dort an. Sie ver\"andert nicht das Produkt (Aumann baut
weiterhin automatisierte Fertigungslinien), sondern den Kundenkreis: Clean
Tech, Aerospace und Life Sciences investieren aus anderen Gr\"unden und zu
anderen Zeitpunkten als die Automobilindustrie. Die technische Kompetenz
bleibt dieselbe, das Konjunkturrisiko verteilt sich auf mehrere Zyklen. F\"ur
ein Unternehmen, dessen Bewertung erkennbar an der Zyklusabh\"angigkeit
h\"angt, ist dies die einzige der drei Optionen, die deren Ursache und nicht
deren Wirkung angeht. Dass es sich nicht um eine blo{\ss}e Ank\"undigung
handelt, zeigt das Segment selbst: Der Auftragseingang von Next Automation
stieg 2025 auf \MioEUR{\AuftragseingangNextAutoZFV} nach
\MioEUR{\AuftragseingangNextAutoZFIV}, w\"ahrend der Segmentumsatz auf
\MioEUR{\UmsatzNextAutoZFV} zur\"uckging und die Segmentmarge sich auf
\Pct{\EbitdaMargeNextAutoZFV} verbesserte (Tabelle~\ref{tab:next-automation}).
Das Segment f\"ullt sich also mit Auftr\"agen, bevor es sie fakturiert.
Einschr\"ankend bleibt die Gr\"o{\ss}enordnung: Gemessen am Segment E-Mobility
mit \MioEUR{\UmsatzEMobilityZFIV} Umsatz ist Next Automation klein. Ein
sp\"urbarer Beitrag zur Risikominderung braucht mehrere Jahre; daf\"ur
sch\"atze ich \MioEUR{\InvestitionsbedarfEigenA} \"uber drei Jahre als
notwendig. Das ist aus der 2023 aktivierten Entwicklungsquote von
\Pct{\FEQuoteZFIII} des Umsatzes hergeleitet, angewandt auf den
Konzernumsatz 2024 \MioEUR{\UmsatzZFIV} und verdreifacht, weil der
Eintritt in drei branchenfremde M\"arkte au{\ss}er Entwicklung auch
Vertrieb, Anwendungstechnik und Zertifizierung verlangt, die nicht
aktivierungsf\"ahig sind.}

\subsubsection{Finanzwirtschaftliche Bewertung}
Tabelle~\ref{tab:einschaetzung-a} fasst die Einsch\"atzung der
Option~A entlang der sechs von der Aufgabenstellung vorgegebenen
Faktoren zusammen. Investitionsbedarf, Kapitalbindung und
Liquidit\"atsbedarf sind eigene Sch\"atzungen aus der
Unternehmensanalyse; Zeithorizont, Risiko und Potenzial nennen
zus\"atzlich die Ankerwerte der Flow-Analyse aus
Abschnitt~\ref{sec:cashflow}, die dort unter den in der
Aufgabenstellung getroffenen Annahmen berechnet sind.

\begin{table}[!htbp]
\centering
\caption{Finanzwirtschaftliche Einsch\"atzung der Option~A.}
\label{tab:einschaetzung-a}
\begin{tabularx}{\linewidth}{@{}p{4.4cm}X@{}}
\toprule
\textbf{Faktor} & \textbf{Einsch\"atzung}\\
\midrule
Investitionsbedarf (F\&E / Markterschlie{\ss}ung)
  & \bk{Mittel: eigene Sch\"atzung \MioEUR{\InvestitionsbedarfEigenA}
    \"uber drei Jahre, also rund \MioEUR{\LiquiditaetsbedarfEigenA}
    j\"ahrlich. Hergeleitet aus der 2023 aktivierten Entwicklungsquote
    von \Pct{\FEQuoteZFIII} des Konzernumsatzes \MioEUR{\UmsatzZFIV},
    verdreifacht f\"ur Vertrieb, Anwendungstechnik und Zertifizierung
    beim Eintritt in drei branchenfremde Zielm\"arkte. Das entspricht
    rund einem Drittel der Nettoliquidit\"at \MioEUR{\NettoliquiditaetZFIV}
    und ist in jeder einzelnen Jahresscheibe aus dem Bestand tragbar.}\\
\addlinespace
Kapitalbindung (Anlagen, CapEx)
  & \bk{Mittel: eigene Sch\"atzung \MioEUR{\KapitalbindungEigenA}, angesetzt
    mit einem Viertel des Investitionsbedarfs als bei Auftragsfertigern
    \"ublichem Verh\"altnis von gebundenem Umlaufverm\"ogen zu
    Projektvolumen. Der Betrag bindet Umlaufverm\"ogen und steht bis zum
    Ende des Vorhabens nicht f\"ur anderes zur Verf\"ugung.}\\
\addlinespace
Liquidit\"atsbedarf (j\"ahrlich)
  & \bk{Gering im Verh\"altnis zum Bestand: eigene Sch\"atzung
    \MioEUR{\LiquiditaetsbedarfEigenA} j\"ahrlich, der Investitionsbedarf
    gleichm\"a{\ss}ig \"uber drei Jahre verteilt. Das entspricht rund
    einem Zehntel der Nettoliquidit\"at \MioEUR{\NettoliquiditaetZFIV}
    und liegt deutlich unter dem, was der Bestand in einem Jahr tr\"agt.}\\
\addlinespace
Zeithorizont bis Return
  & \bk{Lang: erster Ergebnisbeitrag ab \EbitdaStartA, also nach drei
    Investitionsjahren und damit der sp\"ateste der drei Optionen.
    Entsprechend steigt der interne Zinsfu{\ss} (die j\"ahrliche Rendite der
    Investition \"uber ihre gesamte Laufzeit) mit der Laufzeit von
    \Pct{\IrrFuenfA} auf f\"unf auf \Pct{\IrrA} auf zehn Jahre.}\\
\addlinespace
Risiko
  & \bk{Mittel: kein Erwerbs- und kein Integrationsrisiko, aber
    Markteintritt in Branchen ohne breite eigene Referenzbasis. Die
    Investition ist \"uber drei Jahre gestreckt und damit in Stufen
    revidierbar.}\\
\addlinespace
Potenzial
  & \bk{Hoch: mit \MioEUR{\EbitdaZuwachsA} ab \EbitdaStartA\ der gr\"o{\ss}te
    absolute Ergebnisbeitrag der drei Optionen, zugleich als einzige mit
    einer Verringerung der Abh\"angigkeit vom Investitionszyklus der
    Automobilindustrie.}\\
\bottomrule
\end{tabularx}
\end{table}

% Die drei Statement-Ueberschriften folgen unmittelbar auf eine Tabelle und
% tragen darunter nur einen Absatz. Ohne \needspace landet die Ueberschrift
% allein am Seitenfuss - bei 5.2.3 ist genau das passiert.
\needspace{5\baselineskip}
\subsubsection{Statement}
% FAKTENBASIS 5.1.3
%  - Umsetzungsdauer laut Annahme Aufgabe 6: Investition ueber drei Jahre
%    (2026-2028), Ergebnisbeitrag erst ab 2028
%  - Kapitalbindung 15 Mio. EUR belastet das Umlaufvermoegen (current assets)
%  - Liquiditaetsverlauf und Reserve Ende 2028: Tabelle \ref{tab:cf-a}
%  - IRR und Vergleichswerte: Tabelle \ref{tab:vergleich2028},
%    Sensitivitaet Tabelle \ref{tab:irr-sens}
%  - EBITDA-Marge 2024: \EbitdaMargeZFIV %, 2025: \EbitdaMargeZFV %
%  - Segmentmarge Next Automation 2025: \EbitdaMargeNextAutoZFV % - liegt
%    knapp unter der Segmentmarge E-Mobility 2024 (\EbitdaMargeEMobilityZFIV %)
%  - Lehrveranstaltung: Zielkonflikt Rentabilitaet (Eigenkapitalwirkung,
%    verzoegert durch spaeten Ergebnisbeitrag) vs. Liquiditaet
%    (Umlaufvermoegen, sofortige Bindung)
%  - Kerngeschaeft: keine Substitution, sondern Ergaenzung; E-Mobility
%    traegt \AnteilEMobility % des Geschaefts
\bk{Die folgenden Betr\"age und Zeitpunkte sind Ergebnisse der in
Abschnitt~\ref{sec:cashflow} getroffenen Annahmen der Aufgabenstellung,
nicht der eigenen Sch\"atzung aus Abschnitt~\ref{subsec:option-a}.

\textbf{Dauer der Umsetzung.} Option~A ist die langsamste der drei
Optionen. Die Investition von \MioEUR{\InvestitionssummeA} verteilt sich auf
die Jahre 2026 bis 2028, der Ergebnisbeitrag setzt erst \EbitdaStartA\ ein.
Das ist keine Schw\"ache der Ausf\"uhrung, sondern der Sache: Wer in Clean
Tech, Aerospace und Life Sciences als Anlagenbauer auftreten will, braucht
Referenzprojekte, Applikationstechnik und einen Vertrieb, der die
Abnahmekriterien dieser Branchen kennt. Der Vorteil dieser Langsamkeit ist,
dass die Investition in drei Stufen erfolgt und nach jeder Stufe angehalten
oder verkleinert werden kann. Bei einer Akquisition besteht diese
M\"oglichkeit nicht.

\textbf{Rentabilit\"at und Eigenkapitalwirkung.} Der sp\"ate Ergebnisbeitrag
verschiebt die Rentabilit\"atswirkung. Bis \EbitdaStartA\ belasten die
Aufwendungen f\"ur Vertriebsaufbau und Markterschlie{\ss}ung das Ergebnis,
ohne dass ihnen ein Ertrag gegen\"ubersteht; die Eigenkapitalrentabilit\"at
sinkt in dieser Zeit. Erst ab \EbitdaStartA\ kehrt sich das Bild um: Der
Zuwachs von \MioEUR{\EbitdaZuwachsA} ist der gr\"o{\ss}te der drei Optionen
und wirkt \"uber den gesamten weiteren Horizont fort. Der interne Zinsfu{\ss}
zeigt diesen Zeiteffekt deutlich: \Pct{\IrrFuenfA} \"uber f\"unf Jahre,
\Pct{\IrrA} \"uber zehn (Tabelle~\ref{tab:irr-sens}). Wer Option~A \"uber
einen kurzen Horizont beurteilt, beurteilt vor allem ihre Anlaufphase.

\textbf{Liquidit\"at und Umlaufverm\"ogen.} Die Liquidit\"atswirkung ist
zweigeteilt. Die Investition flie{\ss}t gleichm\"a{\ss}ig ab, die
Kapitalbindung von \MioEUR{\KapitalbindungssummeA} kommt in den beiden Jahren
vor dem Ergebnisbeitrag hinzu und erh\"oht das Umlaufverm\"ogen, weil
Anlagenbau den Umsatz vorfinanziert. Beides zusammen zehrt bis Ende 2028 rund
ein Viertel der Ausgangsliquidit\"at auf; es verbleiben
\MioEUR{\LiquiditaetEndeA} (Tabelle~\ref{tab:cf-a}). Damit tritt der in der
Veranstaltung betonte Zielkonflikt in seiner typischen Form auf: Die
Liquidit\"at wird sofort gebunden, die Rentabilit\"at stellt sich erst
sp\"ater ein. Entscheidend ist, dass die Bindung aus eigenen Mitteln
darstellbar bleibt und kein Jahr eine Finanzierungsl\"ucke aufweist.

\textbf{Wirkung auf Kerngesch\"aft und Marge.} Auf das Kerngesch\"aft wirkt
Option~A nicht substituierend, sondern erg\"anzend: Es werden dieselben
Kompetenzen an andere Kunden verkauft. Ein Verdr\"angungseffekt ist allenfalls
bei knappen Konstruktions- und Inbetriebnahmekapazit\"aten zu erwarten. F\"ur
die Marge spricht, dass das Segment Next Automation 2025 mit
\Pct{\EbitdaMargeNextAutoZFV} bereits nahe an der Segmentmarge des
Kerngesch\"afts von \Pct{\EbitdaMargeEMobilityZFIV} lag; die
Diversifikation wird also nicht mit strukturell schlechteren Margen erkauft.
Der ausgewiesene Margensprung auf \Pct{\MargeZweitausendachtundzwanzigA} Ende
2028 ist allerdings zu einem erheblichen Teil ein Effekt der kleinen
Umsatzbasis 2025 und nicht der Option selbst; dazu
Tabelle~\ref{tab:marge-sens}.}

\subsection{Ausbau der Produktionstechnik f\"ur Batterie- und
Brennstoffzellenfertigung (Option B)}
\label{subsec:option-b}

\subsubsection{Relevanz der Strategie (Management-Summary)}
% FAKTENBASIS 5.2.1
%  - Leistungsangebot laut GB 2024, S. 10: Laminier- und Beschichtungsanlagen
%    fuer die Elektroden- und MEA-Fertigung
%  - F&E-Schwerpunkt 2024: Brennstoffzellen- und Elektrolyseurfertigung
%    (GB 2024, S. 13); Batterie-Elektrodenfolien (GB 2024, S. 14)
%  - Akquisition Aumann Lauchheim GmbH, Post-Merger-Integration im
%    Berichtsjahr (GB 2024, S. 6)
%  - Segment E-Mobility 2024: Umsatz \UmsatzEMobilityZFIV Mio. EUR,
%    EBITDA \EbitdaEMobilityZFIII -> \EbitdaEMobilityZFIV Mio. EUR,
%    Marge \EbitdaMargeEMobilityZFIV % (GB 2024, S. 13)
%  - Anteil E-Mobility am Geschaeft: \AnteilEMobility % (GB 2024, S. 10)
%  - Investitionsannahme Aufgabe 6: 30 Mio. EUR ueber drei Jahre,
%    +14 Mio. EUR EBITDA bereits ab 2027 - geringster Kapitalbedarf,
%    frueheste Ergebniswirkung
\bk{Option~B verl\"angert die Wertsch\"opfung dort, wo Aumann bereits arbeitet.
Bisher liefert das Unternehmen vor allem Anlagen f\"ur den elektrischen
Antrieb selbst. Mit Batteriemodul-, Batteriepack- und Cell-to-Pack-Linien
sowie mit Laminier- und Beschichtungsanlagen f\"ur die Elektroden- und
MEA-Fertigung kommen die Fertigungsschritte f\"ur Energiespeicher und
Brennstoffzellen hinzu; Elektrolyse und station\"are Speicher sind
angrenzende Anwendungen desselben Verfahrens.
Wirtschaftlich hat das voraussichtlich das beste Verh\"altnis von
Kapitaleinsatz zu Ertrag: Nach eigener Sch\"atzung verlangt sie mit
\MioEUR{\InvestitionsbedarfEigenB} den geringsten Einsatz, hergeleitet aus
derselben 2023 aktivierten Entwicklungsquote wie Option~A, jedoch nur um
das Eineinhalbfache statt verdreifacht erh\"oht, weil Verfahren, Kunden
und Vertrieb aus dem bestehenden E-Mobility-Gesch\"aft bereits bekannt
sind. Unter den in Abschnitt~\ref{sec:cashflow} getroffenen Annahmen
liefert sie zudem den fr\"uhesten Ergebnisbeitrag bereits ab
\EbitdaStartB\ und erreicht mit \Pct{\IrrB} den h\"ochsten internen
Zinsfu{\ss} des Vergleichs. Der Grund ist einfach: Der
Aufwand entf\"allt \"uberwiegend auf Forschung und Entwicklung und auf den
Markteintritt, nicht auf den Aufbau eines neuen Vertriebs in fremden
Branchen. Die Kompetenz ist zudem im Haus: Die Brennstoffzellen- und
Elektrolyseurfertigung war 2024 ein Schwerpunkt der Entwicklungsarbeit, und
die Akquisition der Aumann Lauchheim GmbH hat das Anlagenspektrum bereits
in diese Richtung erweitert (Abschnitt~\ref{subsec:batterie}).
Die Grenze der Option liegt ebenso klar: Sie verbreitert das Produkt, nicht
den Kundenkreis. Batterie- und Brennstoffzellenfertigung folgt im
Wesentlichen denselben Investitionsentscheidungen der Automobil- und
Energiewirtschaft, an denen Aumann ohnehin h\"angt. Der Anteil des Segments
E-Mobility von rund \Pct{\AnteilEMobility} sinkt dadurch nicht. Option~B
verbessert die Ertragskraft, ohne die Abh\"angigkeit von dieser einen
Kundenbranche zu verringern.}

\subsubsection{Finanzwirtschaftliche Bewertung}
Tabelle~\ref{tab:einschaetzung-b} fasst die Einsch\"atzung der
Option~B entlang der sechs von der Aufgabenstellung vorgegebenen
Faktoren zusammen. Investitionsbedarf, Kapitalbindung und
Liquidit\"atsbedarf sind eigene Sch\"atzungen aus der
Unternehmensanalyse; Zeithorizont, Risiko und Potenzial nennen
zus\"atzlich die Ankerwerte der Flow-Analyse aus
Abschnitt~\ref{sec:cashflow}, die dort unter den in der
Aufgabenstellung getroffenen Annahmen berechnet sind.

\begin{table}[!htbp]
\centering
\caption{Finanzwirtschaftliche Einsch\"atzung der Option~B.}
\label{tab:einschaetzung-b}
\begin{tabularx}{\linewidth}{@{}p{4.4cm}X@{}}
\toprule
\textbf{Faktor} & \textbf{Einsch\"atzung}\\
\midrule
Investitionsbedarf (F\&E)
  & \bk{Gering bis mittel: eigene Sch\"atzung
    \MioEUR{\InvestitionsbedarfEigenB} \"uber drei Jahre und damit
    niedriger als bei Option~A; rund \MioEUR{\LiquiditaetsbedarfEigenB}
    je Jahr. Hergeleitet aus derselben Entwicklungsquote wie Option~A,
    jedoch nur um das Eineinhalbfache statt verdreifacht erh\"oht, weil
    Verfahren, Kunden und Vertrieb bereits bekannt sind.}\\
\addlinespace
Kapitalbindung (Anlagen)
  & \bk{Gering bis mittel: eigene Sch\"atzung \MioEUR{\KapitalbindungEigenB},
    wie bei Option~A mit einem Viertel des Investitionsbedarfs angesetzt.}\\
\addlinespace
Liquidit\"atsbedarf (j\"ahrlich)
  & \bk{Gering: eigene Sch\"atzung \MioEUR{\LiquiditaetsbedarfEigenB}
    j\"ahrlich, der niedrigste Wert der drei Optionen, weil sowohl
    Investitionsbedarf als auch Markterschlie{\ss}ungsaufwand geringer
    ausfallen als bei Option~A.}\\
\addlinespace
Zeithorizont bis Return
  & \bk{Kurz bis mittel: Ergebnisbeitrag ab \EbitdaStartB, also im zweiten
    Jahr. Der interne Zinsfu{\ss} liegt schon auf f\"unf Jahre bei
    \Pct{\IrrFuenfB} und damit \"uber dem Zehnjahreswert der Option~C.}\\
\addlinespace
Risiko
  & \bk{Mittel: technologisch gering, da Verfahren und Kunden bekannt sind;
    marktseitig erh\"oht, weil Hochlauf und F\"ordertempo von Batterie-,
    Brennstoffzellen- und Elektrolyseanwendungen politisch beeinflusst
    sind.}\\
\addlinespace
Potenzial
  & \bk{Hoch bezogen auf die Rendite: \MioEUR{\EbitdaZuwachsB} ab
    \EbitdaStartB\ bei \MioEUR{\KapitalbedarfB} Gesamteinsatz. Begrenzt
    bezogen auf das Risikoprofil, da der Kundenkreis derselbe bleibt.}\\
\bottomrule
\end{tabularx}
\end{table}

% Die drei Statement-Ueberschriften folgen unmittelbar auf eine Tabelle und
% tragen darunter nur einen Absatz. Ohne \needspace landet die Ueberschrift
% allein am Seitenfuss - bei 5.2.3 ist genau das passiert.
\needspace{5\baselineskip}
\subsubsection{Statement}
% FAKTENBASIS 5.2.3
%  - Aufgabenstellung verlangt hier: Kapitalbedarf und Liquiditaetsbelastung
%  - Umsetzungsdauer: Investition 2026-2028, Ergebnisbeitrag bereits ab 2027
%  - Kapitalbedarf 30 Mio. EUR Investition + 15 Mio. EUR Kapitalbindung
%  - Liquiditaetsverlauf und Reserve Ende 2028: Tabelle \ref{tab:cf-b}
%    (tiefster Punkt bereits 2027, Reserve steigt 2028 wieder)
%  - IRR und Vergleichswerte: Tabelle \ref{tab:vergleich2028},
%    Sensitivitaet Tabelle \ref{tab:irr-sens}
%  - Kerngeschaeft: Verbreiterung entlang derselben Wertschoepfungskette,
%    daher hohe technologische Naehe, aber keine Verringerung der
%    Zyklusabhaengigkeit
%  - Marge 2024 \EbitdaMargeZFIV %, 2025 \EbitdaMargeZFV %
\bk{Die folgenden Betr\"age und Zeitpunkte sind Ergebnisse der in
Abschnitt~\ref{sec:cashflow} getroffenen Annahmen der Aufgabenstellung,
nicht der eigenen Sch\"atzung aus Abschnitt~\ref{subsec:option-b}.

\textbf{Dauer der Umsetzung.} Option~B ist die effizienteste der drei
Optionen, gemessen an der Zeit bis zum ersten Ergebnisbeitrag im
Verh\"altnis zum Aufwand. Die Investition l\"auft zwar ebenfalls \"uber drei
Jahre, der Beitrag setzt aber bereits \EbitdaStartB\ und damit ein Jahr
fr\"uher als bei Option~A ein. Der Grund liegt in der Ausgangslage: Verfahren, Kunden und
Vertriebswege sind bekannt, sodass der Markteintritt nicht erst aufgebaut
werden muss.

\textbf{Voraussichtlicher Kapitalbedarf.} Der Gesamtbedarf betr\"agt
\MioEUR{\KapitalbedarfB}, davon \MioEUR{\InvestitionssummeB} Investition und
\MioEUR{\KapitalbindungssummeB} Kapitalbindung. Das ist der niedrigste Bedarf
der drei Optionen und entspricht knapp einem Drittel der
Nettofinanzliquidit\"at von \MioEUR{\NettoliquiditaetZFIV}. Die Finanzierung
aus eigenen Mitteln ist damit unstrittig; eine Fremdfinanzierung w\"are weder
n\"otig noch bei diesem Volumen sinnvoll.

\textbf{Liquidit\"atsbelastung.} Die Belastung ist die geringste des
Vergleichs, und sie ist zeitlich g\"unstig verteilt. Der Aufbau der
Kapitalbindung f\"allt in die Jahre 2026 und 2027, der Ergebnisbeitrag setzt
2027 ein, sodass sich beide Wirkungen bereits im zweiten Jahr teilweise
aufheben. Der Bestand erreicht seinen tiefsten Punkt 2027 und liegt Ende 2028
mit \MioEUR{\LiquiditaetEndeB} wieder \"uber dem Vorjahreswert
(Tabelle~\ref{tab:cf-b}). Die Option finanziert ihre letzte Investitionsrate
also bereits aus dem eigenen Ergebnisbeitrag; das ist der wesentliche
Unterschied zu den Optionen~A und~C, bei denen der Bestand \"uber den gesamten
Betrachtungszeitraum sinkt. Entsprechend f\"allt der interne Zinsfu{\ss} mit
\Pct{\IrrB} \"uber zehn Jahre am h\"ochsten aus, und er ist als einziger
bereits \"uber f\"unf Jahre deutlich positiv (Tabelle~\ref{tab:irr-sens}).

\textbf{Wirkung auf Kerngesch\"aft und Marge.} Die Wirkung auf das
Kerngesch\"aft ist unmittelbar und positiv: Aumann kann demselben Kunden
mehr Fertigungsschritte anbieten und ist im Angebotsvergleich als
Systemlieferant besser gestellt. Das st\"utzt die Marge, weil
Systemgesch\"aft weniger \"uber den Preis einzelner Maschinen entschieden
wird. Diesem Vorteil steht ein Risiko gegen\"uber, das die Kennzahlen nicht
zeigen: Option~B erh\"oht die Konzentration. Aumann w\"urde noch mehr
Ergebnis in derselben Kundenbranche erwirtschaften, deren Investitionszyklus
2024 den Auftragseingang um \PctBetrag{\AuftragseingangDeltaZFIV} einbrechen
lie{\ss}. Die ausgewiesene Marge Ende 2028 von
\Pct{\MargeZweitausendachtundzwanzigB} ist zudem, wie bei allen drei
Optionen, wesentlich vom niedrigen Umsatz des Basisjahres 2025 getrieben
(Tabelle~\ref{tab:marge-sens}). Als finanzwirtschaftliche Einzelentscheidung
ist Option~B die beste des Vergleichs; als alleinige Strategie w\"are sie
unvollst\"andig.}

\subsection{Anorganisches Wachstum durch gezielte Akquisition (Option C)}
\label{subsec:option-c}

\subsubsection{Relevanz der Strategie (Management-Summary)}
% FAKTENBASIS 5.3.1
%  - Nettofinanzliquiditaet 31.12.2024: \NettoliquiditaetZFIV Mio. EUR,
%    Eigenkapital \EigenkapitalZFIV Mio. EUR (GB 2024, S. 5)
%  - Nettofinanzliquiditaet 31.12.2025: \NettoliquiditaetZFV Mio. EUR (GB 2025, S. 4)
%  - Erklaerte Absicht organischen und anorganischen Wachstums (GB 2024, S. 5)
%  - Praezedenzfall: Aumann Lauchheim GmbH, Post-Merger-Integration im
%    Berichtsjahr Gegenstand der Aufsichtsratsberatungen (GB 2024, S. 6)
%  - Reorganisation Aumann Technologies (China) Ltd. (GB 2024, S. 7)
%    -> Integrationsaufwand ist im Konzern belegtes Thema
%  - Investitionsannahme Aufgabe 6: 60 Mio. EUR, davon 45 Mio. EUR Kaufpreis
%    im Vollzugsjahr 2026, Kapitalbindung 20 Mio. EUR, +10 Mio. EUR EBITDA
%    ab 2027 - hoechster Kapitaleinsatz, geringster absoluter Zuwachs
\bk{Option~C kauft Wachstum, statt es aufzubauen. Der Reiz liegt auf der
Hand: Aumann verf\"ugt \"uber eine Nettofinanzliquidit\"at von
\MioEUR{\NettoliquiditaetZFIV}, die im laufenden Gesch\"aft nichts verdient,
und \"uber ein Eigenkapital von \MioEUR{\EigenkapitalZFIV}. Ein passender
Zukauf brächte Technologie, Kunden und Umsatz sofort statt in Jahren, und mit
der Aumann Lauchheim GmbH hat das Unternehmen einen solchen Schritt bereits
vollzogen (Abschnitt~\ref{subsec:batterie}).
Der Kapitalbedarf einer Akquisition l\"asst sich, anders als bei A und~B,
nicht aus der Entwicklungsquote herleiten, da das Zielobjekt offenbleibt.
Nach eigener Sch\"atzung sollte ein einzelner, kaum umkehrbarer
Kapitaleinsatz die H\"alfte der Nettoliquidit\"at \MioEUR{\NettoliquiditaetZFIV}
nicht \"ubersteigen: \MioEUR{\InvestitionsbedarfEigenC} Kaufpreis und
Integration zuz\"uglich \MioEUR{\KapitalbindungEigenC} \"ubernommenes
Umlaufverm\"ogen, zusammen \MioEUR{\LiquiditaetsbedarfEigenC}.
Unter den in Abschnitt~\ref{sec:cashflow} getroffenen Annahmen ist Option~C
die schw\"achste der drei: Sie bindet mit \MioEUR{\KapitalbedarfC} den
h\"ochsten Betrag und liefert mit \MioEUR{\EbitdaZuwachsC} den geringsten
Ergebnisbeitrag. Der interne Zinsfu{\ss} betr\"agt \Pct{\IrrC} \"uber zehn
Jahre, \"uber f\"unf Jahre ist er mit \Pct{\IrrFuenfC} negativ; man zahlt
also das Achtfache des j\"ahrlichen Ergebnisbeitrags, bevor der erste Euro
zur\"uckflie{\ss}t.
Hinzu kommt: Der Kaufpreis von \MioEUR{\KaufpreisC} wird mit dem Vollzug
2026 f\"allig, das Umlaufverm\"ogen im selben Jahr \"ubernommen; der
Liquidit\"atsbestand f\"allt dadurch auf \MioEUR{\LiquiditaetTiefC}
(Tabelle~\ref{tab:cf-c}), und anders als bei A und~B gibt es keinen Punkt
zum Umkehren. Relevant bleibt Option~C daher vor allem als Erg\"anzung, als
kleinerer Zukauf zur Beschleunigung von~A oder~B, nicht als eigenst\"andige
Strategie dieser Gr\"o{\ss}enordnung.}

\subsubsection{Finanzwirtschaftliche Bewertung}
Tabelle~\ref{tab:einschaetzung-c} fasst die Einsch\"atzung der
Option~C entlang der sechs von der Aufgabenstellung vorgegebenen
Faktoren zusammen. Investitionsbedarf, Kapitalbindung und
Liquidit\"atsbedarf sind eigene Sch\"atzungen aus der
Unternehmensanalyse; Zeithorizont, Risiko und Potenzial nennen
zus\"atzlich die Ankerwerte der Flow-Analyse aus
Abschnitt~\ref{sec:cashflow}, die dort unter den in der
Aufgabenstellung getroffenen Annahmen berechnet sind.

\begin{table}[!htbp]
\centering
\caption{Finanzwirtschaftliche Einsch\"atzung der Option~C.}
\label{tab:einschaetzung-c}
\begin{tabularx}{\linewidth}{@{}p{4.4cm}X@{}}
\toprule
\textbf{Faktor} & \textbf{Einsch\"atzung}\\
\midrule
Investitionsbedarf (Kaufpreis \& Integration)
  & \bk{Hoch: eigene Sch\"atzung \MioEUR{\InvestitionsbedarfEigenC}. Ohne
    bekanntes Zielobjekt setze ich als Obergrenze die H\"alfte der
    Nettoliquidit\"at \MioEUR{\NettoliquiditaetZFIV} an, um
    Handlungsf\"ahigkeit f\"ur das laufende Gesch\"aft zu erhalten; davon
    entf\"allt ein Viertel auf die Kapitalbindung.}\\
\addlinespace
Kapitalbindung (Umlaufverm\"ogen erworbenes Gesch\"aft)
  & \bk{Hoch: eigene Sch\"atzung \MioEUR{\KapitalbindungEigenC}, wie bei
    den Optionen~A und~B mit einem Viertel des Investitionsbedarfs
    angesetzt, hier jedoch nicht aufbaubar, sondern mit dem Vollzug
    vollst\"andig zu \"ubernehmen.}\\
\addlinespace
Liquidit\"atsbedarf (j\"ahrlich)
  & \bk{Sehr ungleich verteilt: eigene Sch\"atzung
    \MioEUR{\LiquiditaetsbedarfEigenC} allein im Vollzugsjahr, weil ein
    Kaufpreis nicht in Raten gezahlt wird. Das entspricht rund der
    H\"alfte der Nettoliquidit\"at \MioEUR{\NettoliquiditaetZFIV} und ist
    der h\"ochste Einzeljahresabfluss der drei Optionen.}\\
\addlinespace
Zeithorizont bis Return
  & \bk{Formal kurz (Ergebnisbeitrag ab \EbitdaStartC), wirtschaftlich
    lang: Der interne Zinsfu{\ss} ist \"uber f\"unf Jahre mit
    \Pct{\IrrFuenfC} negativ und erreicht \"uber zehn Jahre nur
    \Pct{\IrrC}.}\\
\addlinespace
Risiko
  & \bk{Hoch: Das Risiko liegt im Kaufpreis und in der Integration, nicht im
    Markt. Es f\"allt in einem einzigen Jahr an und ist danach nicht mehr
    revidierbar; Erfahrungen aus der Integration der Aumann Lauchheim GmbH
    und der Reorganisation in China zeigen, dass dieser Aufwand real ist.}\\
\addlinespace
Potenzial
  & \bk{Begrenzt: \MioEUR{\EbitdaZuwachsC} ab \EbitdaStartC\ ist der
    geringste absolute Beitrag bei h\"ochstem Einsatz. Das Potenzial liegt
    weniger in der Rendite als im Zeitgewinn: Technologie und Marktzugang
    sind sofort vorhanden.}\\
\bottomrule
\end{tabularx}
\end{table}

% Die drei Statement-Ueberschriften folgen unmittelbar auf eine Tabelle und
% tragen darunter nur einen Absatz. Ohne \needspace landet die Ueberschrift
% allein am Seitenfuss - bei 5.2.3 ist genau das passiert.
\needspace{5\baselineskip}
\subsubsection{Statement}
% FAKTENBASIS 5.3.3
%  - Aufgabenstellung verlangt hier: voraussichtlicher Kapitalbedarf
%  - Kapitalbedarf 60 Mio. EUR Investition + 20 Mio. EUR Kapitalbindung,
%    davon 65 Mio. EUR bereits 2026 zahlungswirksam
%  - Tiefster Liquiditaetspunkt und Reserve Ende 2028: Tabelle \ref{tab:cf-c}
%  - IRR ueber zehn Jahre und Sensitivitaet: Tabellen \ref{tab:vergleich2028}
%    und \ref{tab:irr-sens}; ueber fuenf Jahre ist der IRR negativ
%  - Umsetzungsdauer: Vollzug 2026, Ergebnisbeitrag ab 2027 - formal die
%    schnellste Option, aber abhaengig vom Integrationserfolg
%  - Net Debt/EBITDA Ende 2028 laut Tabelle \ref{tab:vergleich2028};
%    negativer Wert bedeutet fortbestehendes Nettoguthaben
%  - Erfahrung aus der Integration der Aumann Lauchheim GmbH (GB 2024, S. 6)
\bk{Die folgenden Betr\"age und Zeitpunkte sind Ergebnisse der in
Abschnitt~\ref{sec:cashflow} getroffenen Annahmen der Aufgabenstellung,
nicht der eigenen Sch\"atzung aus Abschnitt~\ref{subsec:option-c}.

\textbf{Dauer der Umsetzung.} Formal ist Option~C die schnellste: Mit dem
Vollzug 2026 geh\"oren Technologie, Kunden und Umsatz zum Konzern, und der
Ergebnisbeitrag setzt bereits \EbitdaStartC\ ein. Diese Geschwindigkeit
bezieht sich allerdings auf den Erwerb, nicht auf die Wertsch\"opfung. Die
eigentliche Arbeit, die Zusammenf\"uhrung der Prozesse, der Systeme und der
Belegschaften, beginnt erst danach. Der Gesch\"aftsbericht 2024 belegt
diesen Aufwand: Die Post-Merger-Integration der Aumann Lauchheim GmbH und die
Reorganisation der chinesischen Gesellschaft waren im Berichtsjahr Gegenstand
der Aufsichtsratsberatungen.

\textbf{Voraussichtlicher Kapitalbedarf.} Der Bedarf betr\"agt
\MioEUR{\KapitalbedarfC} und ist damit der h\"ochste der drei Optionen. Er
ist nicht nur gr\"o{\ss}er, sondern auch anders strukturiert: W\"ahrend die
Optionen~A und~B ihre Mittel gleichm\"a{\ss}ig \"uber drei Jahre abrufen,
werden bei Option~C der Kaufpreis von \MioEUR{\KaufpreisC} und das
\"ubernommene Umlaufverm\"ogen von \MioEUR{\KapitalbindungssummeC} im selben
Jahr zahlungswirksam. Der Bestand f\"allt dadurch 2026 auf
\MioEUR{\LiquiditaetTiefC} und erholt sich bis Ende 2028 nur auf
\MioEUR{\LiquiditaetEndeC} (Tabelle~\ref{tab:cf-c}). Die Finanzierung bleibt
aus eigenen Mitteln m\"oglich (Net Debt/EBITDA, das Verh\"altnis von Schulden
zum erwirtschafteten Gewinn, bleibt mit \num{\NetDebtEbitdaC} negativ, es
besteht also weiterhin ein Nettoguthaben),
doch der finanzielle Spielraum f\"ur eine zweite Entscheidung w\"are f\"ur
mehrere Jahre aufgebraucht. Genau das ist der eigentliche Preis dieser
Option: nicht die Verschuldung, sondern der Verlust der Handlungsf\"ahigkeit.

\textbf{Rendite und Risiko.} Bei \MioEUR{\KapitalbedarfC} Einsatz f\"ur
\MioEUR{\EbitdaZuwachsC} zus\"atzliches EBITDA ergibt sich ein interner
Zinsfu{\ss} von \Pct{\IrrC} \"uber zehn Jahre; \"uber f\"unf Jahre ist er mit
\Pct{\IrrFuenfC} negativ (Tabelle~\ref{tab:irr-sens}). Damit liegt Option~C
in einer Gr\"o{\ss}enordnung, in der eine Kapitalr\"uckgabe an die Aktion\"are
ernsthaft konkurrenzf\"ahig wird. Hinzu kommt, dass sich das Risiko nicht
staffeln l\"asst: Der Kaufpreis ist mit dem Vollzug versunken, unabh\"angig
davon, ob die unterstellten \MioEUR{\EbitdaZuwachsC} eintreten.

\textbf{Wirkung auf Kerngesch\"aft und Marge.} Sie h\"angt vollst\"andig vom
Zielobjekt ab, das die Aufgabenstellung offenl\"asst. Bei einem
komplement\"aren Automatisierungsunternehmen w\"are der Effekt \"ahnlich dem
der Option~A, bei einem Komponenten- oder Verfahrensunternehmen \"ahnlich dem
der Option~B. Die ausgewiesene Marge von
\Pct{\MargeZweitausendachtundzwanzigC} ist die niedrigste des Vergleichs und
unterliegt demselben Vorbehalt zur Umsatzbasis wie die \"ubrigen
(Tabelle~\ref{tab:marge-sens}). Belastbar ist daher allein die Aussage \"uber
den Kapitaleinsatz; er ist unter den gegebenen Annahmen zu hoch f\"ur
das, was er erbringt.}

\paragraph{Verwendung der Nettoliquidit\"at.}
% FAKTENBASIS 5.3.3 (Zusatz)
%  - Nettofinanzliquiditaet 31.12.2024: \NettoliquiditaetZFIV Mio. EUR,
%    31.12.2025: \NettoliquiditaetZFV Mio. EUR - trotz Ausschuettung und
%    Rueckkauf gestiegen
%  - Dividende je Aktie 2024: \DividendeJeAktieZFIV EUR,
%    Dividendenrendite \DividendenrenditeZFIV %
%  - Rueckkaufangebot ueber bis zu \RueckkaufAktienMax Aktien zu
%    \RueckkaufPreis EUR, anschliessend eingezogen (GB 2024, S. 5)
%  - Aktienzahl \AktienGesamtAlt -> \AktienGesamtMittel -> \AktienGesamt
%  - Kapitalbedarf der Optionen laut Aufgabe 6: A 50, B 45, C 80 Mio. EUR
%    einschliesslich Kapitalbindung
%  - Emissionspreis 2017: \IPOKurs EUR gegenueber Schlusskurs 2025
%    \KursSilvesterZFV EUR
\bk{Diese Frage l\"asst sich nicht allgemein beantworten, sondern nur
\"uber einen Vergleich der Renditen. Kapital an die Aktion\"are
zur\"uckzugeben ist dann richtig, wenn kein eigenes Vorhaben mehr genug
Rendite bringt, um die Kapitalkosten zu decken, also die Mindestrendite, die
Aktion\"are und Kreditgeber f\"ur ihr Geld erwarten. Diese Bedingung ist bei
Aumann nicht erf\"ullt: Die Optionen~A und~B erreichen \Pct{\IrrA}
beziehungsweise \Pct{\IrrB} und liegen damit deutlich \"uber jeder
plausiblen Mindestrendite. Ich w\"urde die Nettoliquidit\"at
daher vorrangig f\"ur Wachstum einsetzen, allerdings f\"ur diese beiden
Optionen und nicht f\"ur eine Akquisition in der unterstellten
Gr\"o{\ss}enordnung, deren \Pct{\IrrC} den Vergleich mit einer
R\"uckgabe kaum besteht.

Die Zahlen st\"utzen diese Reihenfolge auch von der Finanzierungsseite. Die
Optionen~A und~B ben\"otigen zusammen \MioEUR{\KapitalbedarfA} und
\MioEUR{\KapitalbedarfB}, verteilt \"uber drei Jahre; die
Nettofinanzliquidit\"at von \MioEUR{\NettoliquiditaetZFIV} tr\"agt beide
gemeinsam. Vor allem aber ist sie 2025 trotz Dividende und
R\"uckkaufangebot auf \MioEUR{\NettoliquiditaetZFV} gestiegen. Nicht die
Liquidit\"at ist also der Engpass, sondern die Zahl der Vorhaben, die sie
sinnvoll aufnehmen k\"onnen.

Die bisherige R\"uckgabe halte ich gleichwohl f\"ur richtig und w\"urde sie
in ihrem Umfang fortf\"uhren. Sie ist mit einer Dividendenrendite von
\Pct{\DividendenrenditeZFIV} ma{\ss}voll, und der R\"uckkauf zu
\EURje{\RueckkaufPreis} lag weit unter dem Emissionspreis von
\EURje{\IPOKurs}; ein Zukauf eigener Aktien unter dem eigenen historischen
Ausgabepreis ist eine vertretbare Verwendung \"uberschie{\ss}ender Mittel.
Entscheidend ist die Rangfolge: erst die Vorhaben mit belegter Rendite,
dann die R\"uckgabe des Rests, und nicht umgekehrt.}

\subsection{Zusammenfassende Bewertung der drei Optionen}
\label{subsec:uebersicht}

\subsubsection{\"Ubersicht}
Tabelle~\ref{tab:uebersicht-optionen} stellt die drei Optionen entlang
derselben sechs Faktoren nebeneinander und macht damit die Unterschiede
unmittelbar vergleichbar.

\begin{table}[!htbp]
\centering
\footnotesize
\caption{\"Ubersicht der drei Handlungsoptionen. Die Einstufungen folgen einer
einheitlichen Skala (gering / mittel / hoch). Investitionsbedarf,
Kapitalbindung und Liquidit\"atsbedarf sind die eigenen Sch\"atzungen aus den
Tabellen~\ref{tab:einschaetzung-a}, \ref{tab:einschaetzung-b} und
\ref{tab:einschaetzung-c};
Zeithorizont und Chancen nennen die Ankerwerte der Flow-Analyse aus
Abschnitt~\ref{sec:cashflow}.}
\label{tab:uebersicht-optionen}
\begin{tabularx}{\linewidth}{@{}p{1.9cm}XXXXXX@{}}
\toprule
& \textbf{Investitions-} & \textbf{Kapital-} & \textbf{Liquidit\"ats-}
& \textbf{Risiko} & \textbf{Zeit-} & \textbf{Chancen}\\
& \textbf{bedarf} & \textbf{bindung} & \textbf{bedarf (j\"ahrl.)}
& & \textbf{horizont} & \\
\midrule
A: Next Auto\-mation
  & \bk{mittel: \MioEUR{\InvestitionsbedarfEigenA}}
  & \bk{mittel: \MioEUR{\KapitalbindungEigenA}}
  & \bk{gering: \MioEUR{\LiquiditaetsbedarfEigenA} p.\,a.}
  & \bk{mittel: Markt\-eintritt}
  & \bk{lang: ab \EbitdaStartA}
  & \bk{hoch: \MioEUR{\EbitdaZuwachsA}, weniger Zyklus}\\
\addlinespace
B: Batterie / Brenn\-stoff\-zellen
  & \bk{gering: \MioEUR{\InvestitionsbedarfEigenB}}
  & \bk{gering: \MioEUR{\KapitalbindungEigenB}}
  & \bk{gering: \MioEUR{\LiquiditaetsbedarfEigenB} p.\,a.}
  & \bk{mittel: Markt\-hochlauf}
  & \bk{kurz: ab \EbitdaStartB}
  & \bk{hoch: Rendite \Pct{\IrrB}}\\
\addlinespace
C: Akqui\-sition
  & \bk{hoch: \MioEUR{\InvestitionsbedarfEigenC}}
  & \bk{hoch: \MioEUR{\KapitalbindungEigenC}}
  & \bk{hoch: \MioEUR{\LiquiditaetsbedarfEigenC} im Vollzugs\-jahr}
  & \bk{hoch: Kaufpreis, Integration}
  & \bk{kurz formal: ab \EbitdaStartC}
  & \bk{gering: Rendite \Pct{\IrrC}}\\
\bottomrule
\end{tabularx}
\end{table}

\subsubsection{Zusammenfassende Empfehlung}
% FAKTENBASIS 5.4.2
%  - Kumulierte Investitionen: A 35, B 30, C 60 Mio. EUR (Aufgabe 6)
%  - Kapitalbindung: A 15, B 15, C 20 Mio. EUR
%  - EBITDA-Zuwachs: A +16 ab 2028, B +14 ab 2027, C +10 ab 2027
%  - EBITDA-Marge Ende 2028, IRR, Liquiditaetsreserve:
%    Tabelle \ref{tab:vergleich2028}
%  - IRR-Sensitivitaet nach Horizontlaenge: Tabelle \ref{tab:irr-sens}
%  - Ausgangsliquiditaet: \NettoliquiditaetZFIV Mio. EUR - reicht rechnerisch
%    fuer A und B zusammen, nicht fuer alle drei
%  - Auftragseingangsrueckgang 2024: \AuftragseingangDeltaZFIV %,
%    Klumpenrisiko E-Mobility \AnteilEMobility %
%  - Nur Option A adressiert die Zyklusabhaengigkeit selbst
\bk{Empfohlen wird die Verbindung der Optionen~B und~A in dieser Reihenfolge.
Option~B liefert die beste Rendite bei geringstem Einsatz
(\Pct{\IrrB} bei \MioEUR{\KapitalbedarfB}) und tr\"agt ihre letzte
Investitionsrate bereits aus dem eigenen Ergebnis; sie schafft damit den
finanziellen Spielraum f\"ur Option~A. Option~A ist die einzige, die das
eigentliche Problem angeht: die Abh\"angigkeit vom Investitionszyklus einer
Branche, die 2024 einen R\"uckgang des Auftragseingangs um
\PctBetrag{\AuftragseingangDeltaZFIV} verursacht hat. Beide zusammen binden
\MioEUR{\KapitalbedarfB} und \MioEUR{\KapitalbedarfA} und sind aus der
Nettofinanzliquidit\"at von \MioEUR{\NettoliquiditaetZFIV} darstellbar. Von
Option~C in der unterstellten Gr\"o{\ss}enordnung wird abgeraten: h\"ochster
Einsatz, geringster Beitrag, \Pct{\IrrC} \"uber zehn Jahre und kein Punkt zum
Umkehren.}
